%%%%%%%%%%%%%%%%%%%%%%%%%%%%%%%%%%%%%%%%%%%%%%%%%%%%%%%%%%%%%%%%%%%%%%%%%%%%%
%%                                                                         %%
%% \/   \/              Eigene Kommandos definieren                \/   \/ %%
%%                                                                         %%
%%%%%%%%%%%%%%%%%%%%%%%%%%%%%%%%%%%%%%%%%%%%%%%%%%%%%%%%%%%%%%%%%%%%%%%%%%%%%

%%%%%%%%%%%%%%%%%%%%%%%% Zitierern für Anhänge %%%%%%%%%%%%%%%%%%%%%%%%
\zref@newprop*{appsec}{}
\zref@addprop{main}{appsec}

% Definition von \applabel{#1: label}{#2: text}
% von \appsec{#1: text}{#2: label} zur Erzeugung des Labels verwendet
\def\applabel#1#2{%
    \zref@setcurrent{appsec}{#2}%
    \zref@wrapper@immediate{\zref@label{#1}}%
}

% Definition von \appref{#1: label}
% anstelle \ref{#1: label} zum referenzieren von Anhängen verwenden
\def\appref#1{%
    \hyperref[#1]{\zref@extract{#1}{appsec}}%
}
\makeatother

%%%%%%%%%%%%%%%%%%%%%%%% Kaptitel für Anhänge %%%%%%%%%%%%%%%%%%%%%%%%
% Definition von \appsection{#1: text}{#2: label}
% Ersetzt \chapter{#1: text} und \label{#2: label} für Anhänge
\newcommand{\theappsection}[1]{Anhang \Alph{section}:~\protect #1}
\newcommand{\appsection}[2]{
    \addtocounter{section}{1}
    \phantomsection
    \addcontentsline{toc}{section}{\theappsection{#1}}
    \markboth{\theappsection{#1}}{}

    \chapter*{\theappsection{#1}}
    \applabel{#2}{Anhang \Alph{section}}
    \label{#2}
}

%%%%%%%%%%%%%%%%%%%%%%%% In itemize oder enumerate item mit vertikaler linie %%%%%%%%%%%%%%%%%%%%%%%%
\newcommand{\LinedItem}[1]{\item \begin{tabular}[t]{@{}|p{.9847\linewidth}@{}}#1\end{tabular}}
\newcolumntype{P}[1]{>{\centering\arraybackslash}p{#1}}

%%%%%%%%%%%%%%%%%%%%%%%% Unterschtrichene Line am Text %%%%%%%%%%%%%%%%%%%%%%%%
\newcommand{\myuline}[1]{%
    \uline{\phantom{#1}}%
    \llap{\contour{white}{#1}}%
}

%%%%%%%%%%%%%%%%%%%%%%%% Eigene Farbwerte definieren %%%%%%%%%%%%%%%%%%%%%%%%
\definecolor{boxgray}{gray}{0.9}         % Hintergrundfarbe für Zitatboxen
\definecolor{commentgray}{gray}{0.5}     % Grau für Kommentare in Quelltexten
\definecolor{darkgreen}{rgb}{0,.5,0}     % Grün für Strings in Quelltexten

%%%%%%%%%%%%%%%%%%%%%%%% Deutsche Anführungszeichen %%%%%%%%%%%%%%%%%%%%%%%%
% Definition von \gqq{#1: text}: Text in Anführungszeichen
\newcommand{\gqq}[1]{\glqq #1\grqq}

% Definition von \footref{#1: label}
% Verweis auf bereits existierende Fußnoten mittels
\providecommand*{\footref}[1]{
    \begingroup
        \unrestored@protected@xdef\@thefnmark{\ref{#1}}
    \endgroup
\@footnotemark}

%%%%%%%%%%%%%%%%%%%%%%%% Referenzieren mit Seite %%%%%%%%%%%%%%%%%%%%%%%%
% Definition von \mypageref{#1: label}
% Kombination aus \autoref{#1: label} und \pageref{#1: label}
\newcommand{\mypageref}[1]{\autoref{#1} auf Seite \pageref{#1}}

%%%%%%%%%%%%%%%%%%%%%%%% Grauebox für Zitate %%%%%%%%%%%%%%%%%%%%%%%%
% Definition von \myboxquote{#1: text}
% grau hinterlegte Quotation-Umgebung (für Zitate)
\newcommand{\myboxquote}[1]{
    \begin{quotation}
        \colorbox{boxgray}{\parbox{0.78\textwidth}{#1}}
    \end{quotation}
    \vspace*{1mm}
}

%%%%%%%%%%%%%%%%%%%%%%%% Für Bilder (#2 ist das Label) %%%%%%%%%%%%%%%%%%%%%%%%
% Definition \bild{#1 Groesse}{#2 Dateiname}{#3 Bildunterschrift}{#4 Source}
\NewDocumentCommand{\bild}{m m m m o}{%
    \begin{figure}[H]
        \centering
        \includegraphics[width=#1\textwidth, height=#1\textheight, keepaspectratio]{pictures/#2}
        \caption{#3} \label{fig:#2}
        \caption*{\textbf{Source:}%
            \IfNoValueTF{#5}{%
                \cite{#4}%
            }{%
                \cite[p.~#5]{#4}%
            }%
        }
    \end{figure}
}

%%%%%%%%%%%%%%%%%%%%%%%% Für Bilder mit Zitat (#2 ist das Label) %%%%%%%%%%%%%%%%%%%%%%%%
% Definition \bildzitat{#1 Groesse}{#2 Dateiname}{#3 Bildunterschrift}{#4 Seite Zitat}{#5 Zitatkey}
\newcommand{\bildzitat}[5]{
    \begin{figure}[H]
        \centering
        \caption[#3]{#3\protect\footnotemark} \label{fig:#2}
        \includegraphics[width=#1\textwidth, height=#1\textheight, keepaspectratio]{pictures/#2}
    \end{figure}
    \footnotetext{\citet{#5}{#4}}
}

%%%%%%%%%%%%%%%%%%%%%%%% Für Anhang (#2 ist das Label) %%%%%%%%%%%%%%%%%%%%%%%%
% Definition \anhang{#1 Groesse}{#2 Dateiname}{#3 Bildunterschrift}
\newcommand{\anhang}[3]{
    \begin{figure}[H]
        \centering
        \caption{#3} \label{ang:#2}
        \includegraphics[width=#1\textwidth, height=#1\textheight, keepaspectratio]{Anhang/#2}
    \end{figure}
}

%%%%%%%%%%%%%%%%%%%%%%%% Abkürzungen %%%%%%%%%%%%%%%%%%%%%%%%
% Befehl von \abk{#1 Abkürzung}{#2 Erklärung}
\newcommand{\abk}[2]{\nomenclature[]{#1}{#2}}

%%%%%%%%%%%%%%%%%%%%%%%% Zitieren %%%%%%%%%%%%%%%%%%%%%%%%
% Neues Kommando Zitieren \zitat{1 Seite}{2 ID der Quelle}
%\newcommand{\zitat}[2]{\footnote{Vgl. \citet{#2}{#1}}}
%\newcommand{\direktzitat}[2]{\footnote{\citet{#2}{#1}}}

%%%%%%%%%%%%%%%%%%%%%%%% Dezenteres Referenzieren %%%%%%%%%%%%%%%%%%%%%%%%
% Referenziert in Grau und klein \autorefgray{#1 lable}
\newcommand{\autorefgray}[1]{{\footnotesize \color{gray} \textit{\autoref{#1}}}}

%%%%%%%%%%%%%%%%%%%%%%%% Neue Tabellenspalten %%%%%%%%%%%%%%%%%%%%%%%%
\newcolumntype{C}[1]{>{\centering\let\newline\\\arraybackslash\hspace{0pt}}m{#1}}
\newcolumntype{L}[1]{>{\raggedright\let\newline\\\arraybackslash\hspace{0pt}}m{#1}}
\newcolumntype{R}[1]{>{\raggedleft\let\newline\\\arraybackslash\hspace{0pt}}m{#1}}

%%%%%%%%%%%%%%%%%%%%%%%% New z.B. %%%%%%%%%%%%%%%%%%%%%%%%
\newcommand{\forexample}{e.g. }

%%%%%%%%%%%%%%%%%%%%%%%% paragraph with newline %%%%%%%%%%%%%%%%%%%%%%%%
% paragraph mit neuer Zeile \myparagraph{#1 Headline}
\newcommand{\myparagraph}[1]{\paragraph{#1}\mbox{}\\}

%%%%%%%%%%%%%%%%%%%%%%%% subparagraph with newline %%%%%%%%%%%%%%%%%%%%%%%%
% subparagraph mit neuer Zeile \mysubparagraph{#1 Headline}
\newcommand{\mysubparagraph}[1]{\subparagraph{#1}\mbox{}\\}


%%%%%%%%%%%%%%%%%%%%%%%% Glossar %%%%%%%%%%%%%%%%%%%%%%%%
\newglossaryentry{test}{
    name={Test},
    description={Das ist ein Testeintrag}
}
